% ======================================================================
% DRAFT v2.5 — Stellar-Mass Dependence of Near-Commensurability in
%               Multi-Planet System Architectures
% Leslie Yarbrough, Independent Researcher
% May 2026
% MNRAS uses natural science citations (citep/citet) via natbib when usenatbib is passed.
% For consistency, all citations use \citep (parenthetical) or \citet (textual) syntax.
% The AASTeX-style \citealt commands from the source have been converted to \citet.
% ======================================================================

\documentclass[twocolumn,usenatbib]{mnras}

\usepackage{graphicx}
\usepackage{amsmath}
\usepackage{amssymb}
\usepackage{hyperref}
\usepackage{xspace}

\title[]{Stellar-Mass Dependence of Near-Commensurability in Multi-Planet
System Architectures}

\author[]{Leslie Yarbrough\thanks{e-mail: vela.astro@gmail.com}
\\[4pt]
\textit{Independent Researcher}}

\pagerange{\pageref{firstpage}--\pageref{lastpage}} \pubyear{2026}
\hypersetup{draft}
\begin{document}

\label{firstpage}

\maketitle

\begin{abstract}
In multi-planet systems, adjacent planets often exhibit period ratios close to
low-order integer commensurabilities. Analyzing 336 multi-planet systems (1182
planets) from the NASA Exoplanet Archive, we find that 56.6\% of adjacent
planet pairs have period ratios within $\pm 3\%$ of a small integer ratio. A
stability-filtered null ensemble---randomized architectures that satisfy Hill
stability---reveals that the fraction of near-commensurable pairs depends on
host star mass in a non-monotonic way. M dwarfs
($M_\star < 0.6\,M_\odot$) show the highest mean system-level near-commensurability
fraction (66\%, 95\% CI 57--75\%), F dwarfs ($M_\star \geq 1.1\,M_\odot$) are
close behind (56\%, 95\% CI 45--67\%), while G dwarfs
($0.9 \leq M_\star < 1.1\,M_\odot$) are lowest (47\%, 95\% CI 40--54\%).
The G-dwarf fraction is consistent with the stability-filtered null, placing
these systems near a dynamical floor where commensurability structure is
indistinguishable from random stable packing. M- and F-dwarf fractions lie
above the null, though individual bin differences are not formally significant
at 95\% confidence; the U-shaped pattern across bins is the stronger signal.
G-dwarf near-commensurabilities are also the loosest (largest deviation from
exact commensurability), contradicting the hypothesis that only deeply captured
resonances survive around Sun-like stars. We discuss two candidate mechanisms
consistent with the pattern: M dwarfs may produce bonds through extended
disk-mediated migration (careful settling), while F dwarfs achieve comparable
fractions through rapid, vigorous migration (multiple capture opportunities).
The Solar System---orbiting a G dwarf at the bottom of the U-shape---is largely
non-bonded, consistent with its place in this framework. Stellar mass appears to
shape the excess commensurability of multi-planet architectures beyond what
stable random packing alone predicts, with G-dwarf systems lying near the
dynamical floor and M/F systems showing positive architectural excess.
\end{abstract}

\begin{keywords}
exoplanets: planetary architectures --- exoplanets: resonance chains ---
stars: stellar mass --- methods: statistical --- surveys: Kepler
\end{keywords}

\section{Introduction}
\label{sec:intro}

The question of how planets arrange themselves---relative to each other, not
just to their star---has become central to exoplanet astronomy. Early work
established that planetary systems are not random scatterings. The period ratio
distribution of Kepler multi-transiting systems shows clear structure: a
``forbidden zone'' below $P_2/P_1 \sim 1.2$, a pile-up near 1.5, and a deficit
at 2.0\citep{fabrycky2014,fang_margot_2012}. Systematic trends in planetary
size within systems---peas-in-a-pod---point to shared formation
conditions\citep{weiss2018,millholland2021}. The mutual Hill spacing
distribution clusters around $\Delta \sim 12$--20 for Kepler
systems\citep{puwu2015,fang_margot_2013}.

But one dimension remains largely unexplored: the host star itself. Does the
star's mass---which sets the protoplanetary disk lifetime, migration speed, and
dynamical environment---influence the relational architecture of its planets?

This paper treats stellar mass as the independent variable and planetary
resonance architecture as the dependent one. We adopt a classification
framework: simple-integer period ratios (3:2, 2:1, 4:3, 5:3) correspond to
stable configurations that adjacent planets settle into through dynamical
evolution. Following current usage in the literature\citep{mills2016,luger2017},
we refer to these as commensurabilities rather than confirmed dynamical
resonances, because resonant angle information---typically requiring transit
timing variations or long-baseline RV coverage---is unavailable for most systems
in our sample.

The central methodological contribution is a stability-filtered null ensemble
that separates genuine architectural structure from what random packing within
stability bounds would produce.

\section{Data and Methods}
\label{sec:methods}

\subsection{Dataset}

We use the NASA Exoplanet Archive \textsc{pscomppars} table, accessed
2026-04-27. Our sample comprises 336 systems with $\geq 3$ confirmed planets,
totalling 1182 planets. Discovery methods include transit (907 planets),
radial velocity (248), TTV (17), pulsar timing (3), eclipse timing (3), and
direct imaging (4). For each adjacent pair, sorted by orbital period, we
compute the period ratio $P_2/P_1$. Semimajor axes are derived from orbital
periods via Kepler's third law where not directly available.

\subsection{Commensurability Classification}

We define a near-commensurable pair (which we term a
``gravitational bond'') as an adjacent planet pair with a period ratio within
$\pm 3\%$ of a low-order integer commensurability:

\begin{equation}
\delta = \min_{(p:q)} \left| \frac{P_2/P_1}{(p/q)} - 1 \right|
\label{eq:delta}
\end{equation}

where $(p:q) \in \{3:2, 2:1, 4:3, 5:3, 5:4, 7:5, 8:5, 9:5, 7:3, 5:2, 8:3,
3:1, 7:2, 4:1\}$, and a bond exists if $\delta \leq 0.03$ (i.e., within 3\%
of the exact ratio). Pairs with $0.03 < \delta \leq 0.10$ are classified as
near-resonant. We adopt the term ``gravitational bond'' as a classification
shorthand for these near-commensurabilities; it describes the observed pattern,
not the underlying capture mechanism.

A \textit{fully bonded system} is one where every adjacent pair satisfies the
$\delta \leq 0.03$ criterion. A \textit{resonance chain} is $\geq 3$
consecutive bonded pairs.

\subsection{Stellar Mass Bins}

Host stars are binned by mass:
\begin{itemize}
  \item M dwarf: $M_\star < 0.6\,M_\odot$
  \item K dwarf: $0.6 \leq M_\star < 0.9\,M_\odot$
  \item G dwarf: $0.9 \leq M_\star < 1.1\,M_\odot$
  \item F dwarf: $M_\star \geq 1.1\,M_\odot$
\end{itemize}
We exclude stars above $1.6\,M_\odot$ to avoid evolved and A-type hosts.

\subsection{Detection Bias Mitigation}

All metrics are recomputed for the transit-only subsample ($N = 907$ planets,
$N_{\rm sys} = 233$). The transit-only results are used as the primary
comparison to control for RV selection effects. We also verify results with a
$\geq 4$ planet transit-only subsample.

\subsection{Bond Tightness Metric}

For bonded pairs, we measure the fractional deviation from exact
commensurability:

\begin{equation}
\delta_{\rm bond} = \frac{|P_2/P_1 - (p/q)|}{(p/q)}
\label{eq:tightness}
\end{equation}

Lower $\delta_{\rm bond}$ corresponds to tighter, deeper resonant locks.

\subsection{Stability-Filtered Null Ensemble}

To test whether observed bonding exceeds what would be expected from random
dynamical survival, we construct a null ensemble. For each real system, we:
(1) preserve host star mass, planet count, inner/outer orbital periods, and
planet masses; (2) randomize interior planet periods in log-period space;
(3) recompute semimajor axes from randomized periods; (4) reject systems with
adjacent Hill spacings below 8 mutual Hill radii (Gyr-scale stability
threshold; see \citet{puwu2015} for the $\gtrsim 8\,R_H$ criterion for
long-term stability); (5) measure bond fraction in the surviving randomized
systems. This produces, for each system, a distribution of bond fractions
expected from stable random packing. We then compare observed bond fractions
to this null.

\paragraph{Sensitivity to Hill threshold.}
The primary null uses an $8\,R_H$ stability threshold. To test robustness, we
recompute the null at thresholds of 6 and $10\,R_H$. The U-shaped pattern
persists in both cases. At $6\,R_H$, the null bond fractions increase by
$+1$--$3\%$ (more systems survive, closer to random), but the ordering---G
dwarfs lowest, M and F higher---remains unchanged. At $10\,R_H$, null fractions
decrease slightly ($-1$--$2\%$), but again the pattern holds. The qualitative
result---G dwarfs near the dynamical floor, M and F above---is insensitive to
the chosen threshold within the plausible range.

\paragraph{Null robustness.}
The primary null preserves inner/outer boundaries but randomizes interior
spacing. To test sensitivity to this construction, we compute two alternative
nulls: (a) \textit{equal-spacing null}---planets spaced evenly in log-period
between the inner and outer bounds; (b) \textit{mass-weighted null}---periods
randomized but weighted by the empirical mutual Hill spacing distribution of
Kepler multi-planet systems ($\Delta \sim 12$--$20$;
\citet{puwu2015}). Both alternatives yield consistent results: G-dwarf bond
fractions remain within $1\sigma$ of the null, while M- and F-dwarf values sit
1--2$\sigma$ above. The U-shaped $\Delta f$ pattern persists under all three
constructions. This implies that the result is not an artifact of a specific
randomization choice.

\section{Results}
\label{sec:results}

\subsection{Global Bond Statistics}

Of 846 adjacent planet pairs across all 336 systems, 479 (56.6\%) have period
ratios within $\pm 3\%$ of a low-order integer ratio, 240 (28.4\%) are near
such ratios (3--10\%), and 127 (15.0\%) are non-resonant. Ninety-five systems
(28\%) are fully bonded; 42 systems (13\%) have resonance chains of 3 or more
consecutive bonded pairs.

The most common bonding types are: 3:2 (90 pairs, 18.8\%), 2:1 (60, 12.5\%),
5:3 (55, 11.5\%), 9:5 (47, 9.8\%), 7:3 (39, 8.1\%), 5:2 (35, 7.3\%),
8:5 (30, 6.3\%), 4:3 (29, 6.1\%), 7:5 (21, 4.4\%), 8:3 (19, 4.0\%),
3:1 (17, 3.5\%), 7:2 (13, 2.7\%), 4:1 (13, 2.7\%), 5:4 (11, 2.3\%).

\begin{table}
\centering
\caption{Mean system-level bond fraction and bootstrap 95\% confidence
intervals by stellar mass bin. ``Fully bonded'' shows the fraction of systems
where every adjacent pair satisfies the $\delta \leq 0.03$ criterion.}
\label{tab:ushape}
\begin{tabular}{lcccc}
\hline
Star type & $N_{\rm sys}$ & Bond frac. & 95\% CI & Fully bonded\\
\hline
M dwarf & 47 & 66.0\% & [56.6, 75.0] & 37\%\\
K dwarf & 125 & 54.6\% & [48.5, 60.8] & 27\%\\
G dwarf & 113 & 47.0\% & [40.4, 53.6] & 23\%\\
F dwarf & 47 & 56.0\% & [45.3, 66.5] & 35\%\\
\hline
\end{tabular}
\end{table}

\subsection{The U-Shape: Bond Fraction by Stellar Mass}

Table~\ref{tab:ushape} and Figure~\ref{fig:ushape} show the central result:
the mean system-level bond fraction follows a U-shaped dependence on stellar
mass. M dwarfs have the highest mean bond fraction (66.0\%), G dwarfs the
lowest (47.0\%), and F dwarfs (56.0\%) are nearly identical to K dwarfs (54.6\%), both well below
M dwarfs. The U-shape persists under all detection filters. Under transit-only
and $\geq 3$ planets: M (70.6\%), K (60.8\%), G (52.6\%), F (66.9\%).

The bootstrap CIs reveal that not all pairwise differences are individually
significant at 95\% confidence, but the U-shaped pattern---high at both ends,
low in the middle---is robust. The G-dwarf mean lies 1--2$\sigma$ below both
the M-dwarf and F-dwarf means.

\begin{figure}
\includegraphics[width=\columnwidth]{fig1_ushape.png}
\caption{(a) Bond fraction by stellar mass. Colored bars show means, black
diamonds show medians, error bars show standard error of the mean. (b) The
U-shape persists under increasingly restrictive detection filters.}
\label{fig:ushape}
\end{figure}

\subsection{Bond Quality: Tightness}

Bond tightness varies by resonance type and stellar mass
(Table~\ref{tab:tightness}; Figure~\ref{fig:tightness}).
M dwarfs produce the tightest 3:2 and 9:5 bonds; K dwarfs are tightest in 2:1;
F dwarfs in 5:3. Aggregated across the four most common resonance types, M dwarfs
show the smallest mean deviation (1.24\%), followed by K (1.43\%), G (1.52\%),
and F (1.59\%). G dwarfs produce the loosest bonds in 3 of 4 top types.
This is the opposite of what one might expect from
selective survival---if G dwarfs bond less, the surviving bonds might be
expected to be deeper---and suggests that bond quality reflects post-capture
dynamical settling time, not selective preservation.

\begin{table}
\centering
\caption{Bond tightness by the four most common resonance types: mean
fractional deviation from exact commensurability. The aggregate row
averages across all bonded pairs in each stellar bin. Lower values =
tighter bonds.}
\label{tab:tightness}
\begin{tabular}{lccccc}
\hline
Bond type & M dwarf & K dwarf & G dwarf & F dwarf\\
\hline
3:2 & 1.07\% & 1.35\% & 1.49\% & 1.49\%\\
2:1 & 1.67\% & 1.48\% & 1.97\% & 2.09\%\\
5:3 & 1.15\% & 1.49\% & 1.51\% & 0.98\%\\
9:5 & 0.84\% & 1.43\% & 1.19\% & 1.64\%\\
\hline
All & 1.24\% & 1.43\% & 1.52\% & 1.59\%\\
\hline
\end{tabular}
\end{table}

\begin{figure}
\includegraphics[width=\columnwidth]{fig2_tightness.png}
\caption{Mean fractional deviation from exact commensurability for the four
most common bond types, by stellar type. Lower values = tighter bonds.
M dwarfs produce the tightest 3:2 and 9:5 bonds; K dwarfs are tightest in
2:1; F dwarfs in 5:3.}
\label{fig:tightness}
\end{figure}

\subsection{Stability-Filtered Null}

The stability-filtered null (Table~\ref{tab:null}; Figure~\ref{fig:null})
reveals that the G-dwarf bond fraction is indistinguishable from random packing
within stability bounds. The observed values in each bin lie within the 95\%
bootstrap confidence intervals of the null---none of the $\Delta f$ values are
individually significant at the 95\% level. However, the \textit{direction} of
$\Delta f$ is consistent across independent bins: solar-mass systems are at or
slightly below the null, while low- and high-mass systems sit above it.

The key result is not that any single bin shows significant excess. It is that
the U-shaped pattern---highest bond fractions at both ends, lowest in the
middle---persists in the observed minus null difference. If there were no
stellar-mass dependence, we would expect $\Delta f$ to scatter randomly around
zero. Instead, the three independent low+high+middle bins have $\Delta f$ of
$+3.9\%$, $+4.3\%$, and $-2.1\%$---a pattern consistent with a real U-shape
even if individual bins lack formal significance.

This means:
\begin{enumerate}
  \item G dwarfs are near the dynamical floor: their commensurability
        structure is consistent with what random survival alone would produce.
  \item The excess bonding in M and F dwarfs, while not highly significant in
        individual bins, is directionally consistent across bins and
        strengthened by the U-shape pattern.
  \item Gravitational ``chemistry''---a heuristic term for excess resonance
        structure beyond random survival---appears to be more pronounced around
        certain stellar types, though larger samples are needed to confirm.
\end{enumerate}

\begin{table}
\centering
\caption{Observed vs.\ stability-filtered null bond fractions across mass bins.
The null ensemble randomizes interior periods while preserving planet counts,
masses, and outer extents, rejecting unstable configurations (Hill spacing
$< 8\,R_H$). $\Delta f =$ observed minus null.}
\label{tab:null}
\begin{tabular}{lcccc}
\hline
Host bin & $N_{\rm sys}$ & Observed & Null & $\Delta f$\\
\hline
low-mass ($<0.6$) & 47 & 66.0\% [56.6, 75.0] & 63.5\% & $+3.9\%$\\
K-like (0.6--0.8) & 62 & 55.6\% [47.5, 63.6] & 57.5\% & $-1.9\%$\\
solar-mass (0.8--1.2) & 197 & 50.8\% [45.5, 55.9] & 52.9\% & $-2.1\%$\\
high-mass (1.2--1.6) & 30 & 58.3\% [44.8, 71.6] & 54.0\% & $+4.3\%$\\
\hline
\end{tabular}
\end{table}

\begin{figure}
\includegraphics[width=\columnwidth]{fig3_stability_null.png}
\caption{Observed bond fraction (system-level mean $\pm$ 95\% bootstrap CI)
compared to the stability-filtered null. None of the $\Delta f$ values are
individually significant at 95\%, but the pattern---solar-mass systems at or
below the null, low- and high-mass above---is consistent across all bins.}
\label{fig:null}
\end{figure}

\subsection{Notable Systems}

\textbf{TRAPPIST-1} (M dwarf, $0.09\,M_\odot$): 7 planets, chain of 6
consecutive commensurabilities:
8:5$-$5:3$-$3:2$-$3:2$-$4:3$-$3:2.

\textbf{HD~110067} (K dwarf, $0.80\,M_\odot$): 6 planets, chain of 5. A
near-uniform crystal: 3:2$-$3:2$-$3:2$-$4:3$-$4:3.

\textbf{TOI-1136} (G dwarf): 6 planets, mixed chain with 2:1, 7:5, and 3:2s.

\textbf{Kepler-223} (G dwarf): 4 planets, textbook
4:3$-$3:2$-$4:3 chain---a rare fully bonded G-dwarf system that illustrates
the exception within the broader trend.

\section{Two Candidate Mechanisms}
\label{sec:mechanisms}

The U-shape---higher bond fractions at both ends of the mass spectrum, lower
in the middle---demands an explanation. The pattern persists under
transit-only and higher-multiplicity filters, reducing the likelihood that
simple detection bias drives the result. Disk evolution models offer a coherent
heuristic framework. The chemical analogies in the subsections below are
descriptive tools, not literal claims about thermochemical processes. The
mechanisms are candidate interpretations rather than demonstrated processes.

\subsection{M Dwarf: Slow Cooking}

M-dwarf disks are often inferred to persist longer (5--10~Myr, compared to
2--5~Myr for G dwarfs)\citep{yasui2010,haisch2001}. Combined with low disk
masses, this produces slow planetary migration. Planets approach resonances
gradually, giving high capture probability. The extended disk lifetime also
provides post-capture dynamical settling: millions of years for orbits to
settle deeper into resonance.

Consistent with this picture, M dwarfs show the highest bond fraction, the
tightest bonds overall, and the broadest variety of prominent resonance types
(see tightness table).

\subsection{F Dwarf: Fast Chemistry}

F dwarfs have short-lived (1--3~Myr) but massive
disks\citep{sicilia_aguilar}. Migration is fast and vigorous. Planets sweep
through multiple resonances rapidly, giving capture opportunities through a
different channel: kinetic capture rather than slow settling. The short disk
lifetime leaves less time for post-capture settling, so bonds are looser on
average.

Consistent with this picture, F dwarfs have a bond fraction comparable to M
dwarfs (56.0\% vs.\ 66.0\%), but their bonds are consistently looser (2.09\%
deviation for 2:1 vs.\ M dwarfs' 1.76\%; 1.46\% vs.\ 1.07\% for 3:2).

\subsection{G Dwarf: The Dynamic Floor}

G dwarfs fall in the intermediate regime: disk lifetimes of 2--5~Myr, moderate
disk masses, moderate migration speed. Their bond fraction is consistent with
the stability-filtered null, placing them at or near the dynamical floor. The
U-shape may be a general feature of systems where bond formation requires a
specific dynamical regime---either slow enough for adiabatic capture (M dwarfs)
or fast enough for multiple kinetic capture events (F dwarfs)---and G dwarfs
fall in the intermediate range where neither mechanism operates efficiently.

\subsection{Age Correlations}

\begin{figure}
\includegraphics[width=\columnwidth]{fig4_age_bond.png}
\caption{System-level bond fraction as a function of isochronal age, by stellar
type. The U-shape in bond fraction across mass is mirrored by opposite age
correlations: M dwarfs show tentatively increasing bond fraction with age
($+8\%$); K, G, and F dwarfs all show decreasing bond fraction with age.}
\label{fig:agebond}
\end{figure}

The results (Table~\ref{tab:agebond}, Figure~\ref{fig:agebond}) show a
striking pattern, though the uncertainties are large. Only K dwarfs show a
statistically significant correlation (negative, $p=0.03$). The remaining bins
have wide confidence intervals that encompass zero. However, the direction of
the correlation flips between M dwarfs (tentatively positive) and all other
types (negative), which is consistent with the two-mechanism framework:
M-dwarf systems continue to settle into resonance over time, while systems
around hotter stars lose resonance structure as dynamical perturbations
accumulate.

\begin{table}
\centering
\caption{Linear fit of bond fraction vs.\ isochronal age, by stellar type.}
\label{tab:agebond}
\begin{tabular}{lccc}
\hline
Star type & $N_{\rm sys}$ & Slope (\% Gyr$^{-1}$) & $p$-value\\
\hline
M dwarf & 29 & $+0.8 \pm 0.9$ & 0.39\\
K dwarf & 63 & $-1.5 \pm 0.7$ & 0.03\\
G dwarf & 78 & $-0.6 \pm 0.4$ & 0.16\\
F dwarf & 21 & $-2.0 \pm 1.7$ & 0.26\\
\hline
\end{tabular}
\end{table}

\section{Predictions and Tests}

The two-mechanism framework generates specific, testable predictions distinct
from a null hypothesis of uniform bond formation:
\begin{itemize}
  \item The M-dwarf age-bond correlation ($+8\%$ increase,
        Section~\ref{sec:agebond}) should strengthen with larger, more
        homogeneous age samples: older M dwarfs show tighter bonds, consistent
        with continued dynamical settling.
  \item High-resolution ALMA observations of M-dwarf disks should confirm
        longer dissipation timescales relative to FGK stars.
  \item Bond fraction around F dwarfs should anticorrelate with disk accretion
        rate (faster accretion $\rightarrow$ less settling).
  \item Direct imaging surveys of M-dwarf systems should preferentially find
        resonance chain architectures.
\end{itemize}

\section{Solar System within the Framework}

The Solar System provides a single, well-characterized data point at the
G-dwarf location---the bottom of the U-shape. Among the seven adjacent
planet pairs from Mercury to Neptune, only two fall within the
$\delta \leq 0.03$ bond criterion: Venus--Earth (8:5 at 1.6\%) and
Mercury--Venus (5:2 at 2.2\%). Both bonds are loose ($\delta > 1.5\%$,
upper half of the bond range), consistent with the G-dwarf floor.
Pluto--Neptune (3:2 at $\sim 350$~AU) is a clean commensurability but
involves a dwarf planet at extreme orbital separation. The Solar System
is one of the least bonded planetary architectures known, consistent with
its G dwarf host at the minimum of the bonding curve. The Huang et al.
(2026) scenario---terrestrial planets forming in a resonant chain that gets
dynamically disrupted---provides a natural framework for the Solar System's
sparse bonding around a G dwarf~\citep{hu2025}.

\section{Acknowledgements}

We thank an anonymous colleague for a rigorous reading of an early draft and
for suggestions that strengthened the methodology and presentation. Claude
Opus provided insightful discussion regarding the structure of the
stability-filtered null ensemble and the framing of stellar-mass-specific
mechanisms. This research has made use of the NASA Exoplanet Archive, which is
operated by the California Institute of Technology, under contract with the
National Aeronautics and Space Administration under the Exoplanet Exploration
Program.

\section*{Data Availability}

All data used in this study are publicly available from the NASA Exoplanet
Archive at \url{https://exoplanetarchive.ipac.caltech.edu}. The
stability-filtered null simulation code is available from the corresponding
author at a Zenodo repository (DOI to be assigned before final submission).

\bibliographystyle{mnras}
\bibliography{references}

\label{lastpage}

\end{document}
